\subsubsection{One-Nearest-Neighbour Classification}\label{sec:one-nearest-neighbour-classification}

In this section, we will make a very small change to the image-retrieval pipeline we have built in the previous section. \hl{Instead of only retrieving the closest image, use its label as the prediction for the query image.} We will implement a \textbf{1-NN classification in the latent space} by extracting the query feature vector, extracting all training-set feature vectors, computing distances, sorting them, and taking the label associated with the nearest training image.
\begin{lstlisting}[language=Python]
# feed the test image to the FEN
image_features = fen.predict(test_image)

# feed the entire training set to the FEN
features = fen.predict(
    X_train_val,
    batch_size=512,
    verbose=0
)

# compute L1-like distances
distances = np.mean(
    np.abs(features - image_features),
    axis=-1
)

# sort by increasing distance
sortedDistances = distances.argsort()

# reorder images and labels
ordered_images = X_train_val[sortedDistances]
ordered_labels = y_train_val[sortedDistances]

# closest sample is ordered_images[0]
\end{lstlisting}
Let's see step by step how this works:
\begin{enumerate}
    \item \textbf{Extract the query feature vector}:
    \begin{lstlisting}[language=Python]
# feed the test image to the FEN
image_features = fen.predict(test_image)
\end{lstlisting}
    Given a query image $I_q$:
    \begin{equation*}
        \mathbf{z}_q = \text{FEN}(I_q)
    \end{equation*}
    The original CNN classifier is not needed here. We only use the \textbf{Feature Extraction Network (FEN)}.
    
    
    \item \textbf{Extract the training representations}:
    \begin{lstlisting}[language=Python]
# feed the entire training set to the FEN
features = fen.predict(
    X_train_val,
    batch_size=512,
    verbose=0
)
\end{lstlisting}
    For each training image $I_i$:
    \begin{equation*}
        \mathbf{z}_i = \text{FEN}(I_i)
    \end{equation*}
    So training set becomes a labelled feature database:
    \begin{equation*}
        \mathcal{D} = \left\{
            \left(
                \mathbf{z}_i , y_i
            \right)
        \right\}_{i=1}^{N}
    \end{equation*}
    In the code snippet above, we use a batch size of 512 to speed up the feature extraction process. The batch indicates how many images are processed in parallel by the FEN.


    \item \textbf{Compute the distance to every training feature}:
    \begin{lstlisting}[language=Python]
# compute L1-like distances
distances = np.mean(
    np.abs(features - image_features),
    axis=-1
)
\end{lstlisting}
    We compute the distance between the query feature vector and every training feature vector. In this case, we use a \textbf{L1-like distance}:
    \begin{equation*}
        d(\mathbf{z}_i, \mathbf{z}_q) = \frac{1}{D} \sum_{j=1}^{D} \left| z_{i,j} - z_{q,j} \right|
    \end{equation*}
    Where $D$ is the dimension of the latent space. The result is a vector of distances:
    \begin{equation*}
        \mathbf{d} = [d(\mathbf{z}_q, \mathbf{z}_1), d(\mathbf{z}_q, \mathbf{z}_2), ..., d(\mathbf{z}_q, \mathbf{z}_N)]
    \end{equation*}
    Then all distances are sorted from smallest to largest:
    \begin{lstlisting}[language=Python]
# sort by increasing distance
sortedDistances = distances.argsort()
\end{lstlisting}


    \item \textbf{Reorder the training images and labels}:
    \begin{lstlisting}[language=Python]
# reorder images and labels
ordered_images = X_train_val[sortedDistances]
ordered_labels = y_train_val[sortedDistances]
\end{lstlisting}
    The training images and labels are reordered according to the sorted distances. The closest training image is now the first element of the ordered list:
    \begin{equation*}
        I_{closest} = \texttt{ordered\_images}[0]
    \end{equation*}
    \begin{equation*}
        y_{closest} = \texttt{ordered\_labels}[0]
    \end{equation*}


    \item \textbf{Copy the label of the closest sample}. Once the nearest representation has been found $\mathbf{z}_{i^*}$, where:
    \begin{equation*}
        i^* = \argmin_{i} d\left(\mathbf{z}_q, \mathbf{z}_i\right)
    \end{equation*}
    The predicted label is simply the label of the closest training image:
    \begin{equation*}
        \hat{y}_q = y_{i^*}
    \end{equation*}
    So the classifier is simply a \textbf{1-NN classifier in the latent space}:
    \begin{equation*}
        \hat{y}_q = y_{\argmin_{i} d\left(\mathbf{z}_q, \mathbf{z}_i\right)}
    \end{equation*}
\end{enumerate}
In contrast to the image-retrieval pipeline, we do not need to display the retrieved images. We only need to return the label of the closest training image:
\begin{equation*}
    \text{query} \to \text{nearest latent vector} \to \text{take its label}
\end{equation*}
In summary, \hl{the CNN is valuable not only as a classifier, but also as a learned feature extractor whose latent space can support retrieval and even a simple nearest-neighbour classifier.}

\highspace
\begin{takeawaysbox}
    \textbf{1-NN in latent space replaces the CNN's original classifier with a nearest-neighbour rule.} The FEN is retained to produce meaningful feature vectors; the query is classified with the label of the training image whose latent representation is closest.
\end{takeawaysbox}